\documentclass[11pt,a4paper]{article}

% --- Packages ---
\usepackage[margin=1in]{geometry}
\usepackage{graphicx}
\usepackage{booktabs}
\usepackage{hyperref}
\usepackage{amsmath}
\usepackage[T1]{fontenc}
\usepackage{authblk}

\title{Auditing the OPAS Orbital Proximity Alert System's Safe-Window Scanner}
\author[1]{Nadim Baboun}
\affil[1]{Al-Quds University}
\date{\today}

\begin{document}

\maketitle

\begin{abstract}
The Orbital Proximity Alert System (OPAS) is a web-based collision avoidance tool that estimates safe launch windows from two-line element (TLE) sets and the SGP4 propagation model. This study audits its \texttt{/safe-windows} endpoint against a dense, independently validated reference calculation. Every window the original scanner reported as safe in three crowded low Earth orbit shells (ISS, Starlink, and sun-synchronous) was, per the reference, unsafe for its entire duration, while a geostationary transfer control case was correct. The failure is traced to insufficient discrete trajectory sampling relative to object density. A corrected scanner---dense 600-point sampling, an age-scaled 10--15~km screening radius, and one-minute temporal verification---returns one window each for ISS (9.6\% false-safe) and Starlink (21.1\% false-safe), with all ten residual close approaches traced to the corrected scanner's finite waypoint spacing, and no window for SSO. Direct enumeration of clear launch-time intervals in the dense reference, independent of the scanner's own window-forming algorithm, confirms that no interval of at least 15 minutes exists for Starlink or SSO at the corrected radius --- so the reported Starlink window is itself false-safe, and SSO's absence of any window is correct --- while showing that the same algorithm can also miss genuinely clear ISS windows that do exist. We compare these choices with documented NASA conjunction-assessment practice and propose anisotropic, continuous swept-volume screening as future work.
\end{abstract}

\section{Introduction}

The Orbital Proximity Alert System (OPAS) is a lightweight web tool for assessing
collision risk during satellite and rocket launches. Given a launch site, target
orbit, and inclination, its \texttt{/safe-windows} endpoint scans forward in time
and returns time intervals during which, according to the tool, the planned
trajectory stays clear of tracked orbital debris and satellites\footnote{Source
code: \url{https://github.com/nadeemtsf/OPAS}.}. The system ingests two-line
element (TLE) sets from Space-Track, propagates object positions using the SGP4
model via the Skyfield library, and screens a simplified circular-orbit vehicle
trajectory against a catalog of tracked objects near the target altitude.

OPAS was developed through rapid, AI-assisted prototyping rather than a formal
aerospace engineering process. Its core output --- a claimed ``safe launch
window'' --- is precisely the kind of claim that should not be trusted without
independent verification, particularly given that the tool's internal screening
logic relies on coarse temporal and spatial sampling to keep response times
practical for a web API. This report presents an empirical audit of that
screening logic: whether the windows OPAS reports as safe are, in fact, free of
close approaches, and if not, why.

To answer this, we built a dense, independently-implemented reference
calculation that estimates the true minimum distance between the launch
vehicle and every tracked object at one-second resolution, using the same
underlying orbital data and propagation model as OPAS itself, so that any
disagreement between the two can be attributed to sampling density rather than
a difference in orbital mechanics. We then validated OPAS's actual returned
safe windows against this reference across five representative launch
scenarios spanning low Earth orbit (LEO), a sun-synchronous orbit, and a
geostationary transfer orbit.

The remainder of this report is organized as follows. Section~\ref{sec:background}
reviews relevant literature on TLE/SGP4 propagation accuracy and documented
conjunction-screening practice. Section~\ref{sec:methodology} describes the
frozen data snapshot, reference implementation, and validation harness.
Section~\ref{sec:results} presents the audit's findings: the original
scanner's false-safe rate, the design flaws responsible for it, an engineered
fix, and a direct, case-by-case analysis of the fix's remaining limitation.
Section~\ref{sec:discussion} interprets these results against documented
real-world practice, Section~\ref{sec:limitations} states the study's scope
boundaries, and Section~\ref{sec:conclusion} summarizes findings and proposes
future work.

\section{Background}
\label{sec:background}

\subsection{TLE/SGP4 propagation accuracy}

OPAS propagates object positions from two-line element (TLE) sets using the
SGP4 model. TLE/SGP4 accuracy degrades with time since the TLE's epoch: a
commonly cited heuristic places initial error around 1~km at epoch, growing by
an additional 1--3~km per day thereafter for low Earth orbit (LEO)
objects~\cite{jaxsgp4}. An empirical comparison against a numerical
integrator found SGP4-propagated position error for a real LEO object stayed
within 1~km for roughly the first week from epoch, with error growing
super-linearly afterward and dominated by the in-track (along-orbit)
component rather than the radial or cross-track directions~\cite{afit_sgp4}.
This distinction --- error growing much faster along-track than radially ---
is directly relevant to how a screening radius should be shaped, discussed
further in Section~\ref{sec:discussion}.

\subsection{Documented conjunction-screening practice}

NASA's Conjunction Assessment and Collision Avoidance Best Practices Handbook
states that screening volumes are assigned per protected asset and orbit
regime rather than fixed~\cite{nasa_handbook}. One documented coarse-screening
pass, used across a six-month operational dataset, applied a volume of
approximately 50~km (radial) by 250~km (in-track) by 250~km (cross-track),
looking several days ahead~\cite{nasa_screening}. A separate documented
operational example, for a near-Earth-orbit protected asset, instead uses a
much tighter 0.5~km~$\times$~17~km~$\times$~20~km ellipsoid, reported
independently in both sources~\cite{nasa_screening,nasa_handbook}. A related
NASA presentation states that TLE-based tracking accuracy --- on the order
of 1--2~km within a few days of epoch --- is not considered sufficient for
official conjunction assessment by the 18th Space Control Squadron, and that
operators requiring rigorous collision avoidance instead submit their own
ephemerides for direct screening~\cite{nasa_cospar}.

Two points from this practice inform the present study. First, real
screening volumes are anisotropic ellipsoids, not spheres: even the
tightest documented example (0.5~km radial) uses a substantially larger
in-track and cross-track extent (17--20~km), reflecting the
in-track-dominated error behavior discussed above. Second, screening
volumes are assigned per mission rather than fixed, and even the largest
documented example (50~km radial) was used only as an initial coarse
filter within one study, not applied universally as a final safety
threshold --- a distinction relevant to evaluating OPAS's single, final,
universally-applied 50~km spherical radius (Section~\ref{sec:results}).

\section{Methodology}
\label{sec:methodology}

\subsection{Frozen data snapshot}

To ensure every experiment in this study ran against identical data, the
live \texttt{opas\_db.debris} MongoDB collection was exported once to a
static, checksummed snapshot (32{,}517 objects, SHA-256 verified before every
subsequent script run). TLE age was computed relative to this snapshot's
export time throughout, rather than the wall clock, so that results are
reproducible independent of when the analysis scripts are actually executed.
All measurements below ran against this single snapshot and the unmodified
OPAS source at git commit \texttt{a89f2e72}\footnote{Python 3.13.14, NumPy
2.4.4, Skyfield 1.54; OPAS's optional native C++ distance-check extension
was active for all runs.}.

\subsection{Baseline measurement}

The unmodified \texttt{scan\_windows} function was run against five launch
presets spanning a range of orbital regimes: an ISS resupply orbit (420~km,
51.6\textdegree{} inclination), a Starlink deployment orbit (550~km,
53.0\textdegree{}), a sun-synchronous orbit (705~km, 98.2\textdegree{}), a
polar orbit (800~km, 90\textdegree{}), and a geostationary transfer orbit
(35{,}786~km, 6.0\textdegree{}), each scanned over a 6-hour horizon. This
produced, for each preset, the exact set of ``safe windows'' OPAS itself would
report to a user.

\subsection{Dense reference implementation}

To evaluate whether those reported windows were actually free of close
approaches, we built an independent reference calculation using the same
orbital data, propagation model (SGP4 via Skyfield), and simplified
circular-orbit vehicle trajectory as OPAS, differing only in sampling
density and closest-approach interpolation: object and vehicle positions were
sampled at approximately 1-second resolution (versus OPAS's coarser internal
sampling, detailed in Section~\ref{sec:results}), and the true closest-approach
distance between sample points was recovered via parabolic interpolation of
the three nearest samples --- exact for locally straight-line relative
motion. Because both OPAS and the reference share the same underlying
orbital mechanics, any disagreement between them is attributable to how
finely closest approach is resolved rather than a difference in physics.

The reference implementation was validated before use: a self-test using
two planted objects with known closest-approach geometry recovered the
correct minimum distance to within 0.02~km after parabolic refinement, and
an independent comparison of 1-second versus 2-second sampling across
38{,}000 object-launch pairs showed a median disagreement of 0.5~m with zero
classification flips at any of four candidate screening radii (5, 10, 25, and
50~km), confirming that 1-second sampling is dense enough to serve as ground
truth at the distances relevant to this study.

\subsection{Validating OPAS's reported windows}

For each preset, every window OPAS's baseline run reported as safe was
checked minute-by-minute against the dense reference at OPAS's own
per-object screening radius, yielding a \emph{false-safe rate}: the
fraction of a reported window's duration during which the reference showed
at least one tracked object inside that same radius. A parameter sweep
across OPAS's internal sampling settings (coarse step, fine step, and
inner-orbit sample count) was additionally run to test whether the
false-safe rate was attributable to sampling density specifically, as
opposed to some other property of the scanning algorithm. The parameterized
scanner used for this sweep was confirmed to reproduce \texttt{scan\_windows}'s
baseline output exactly (identical windows and evaluation counts) for the
three presets it was run against (ISS, Starlink, SSO) before its results were
used.

\subsection{Engineered fix and validation}

Based on the findings in Section~\ref{sec:results}, a corrected scanner was
implemented with four changes: (i) trajectory sampling density increased
from roughly a dozen to 600 points per orbit, (ii) the per-object screening
radius rescaled from a flat 50~km with an exponential TLE-age penalty to a
10~km base with a capped linear age penalty (Section~\ref{sec:discussion}),
(iii) temporal verification tightened from 2-minute to 1-minute resolution,
and (iv) returned windows clamped to the requested scan range. This
corrected scanner's output was validated against the dense reference using
the exact per-object radius formula it applies internally --- rather than an
approximate flat radius --- so that the reported false-safe rate reflects the
scanner's own actual decision criterion precisely. Where a non-zero
false-safe rate remained, its cause was examined directly: for each
reference-confirmed close approach inside a corrected-scanner ``safe''
window, the scanner's own sampled positions at that exact launch time were
recomputed to confirm whether its own math, run on its own inputs, indeed
reported no threat. Separately, every contiguous stretch of reference-confirmed
clear one-minute launch times was enumerated directly from the reference data,
independent of the scanner's own window-forming algorithm, to test whether a
genuinely clear window of useful length exists regardless of whether the
scanner itself located it.

\section{Results}
\label{sec:results}

\subsection{Baseline: reported safe windows were false-safe in crowded shells}

Table~\ref{tab:baseline} summarizes every window OPAS's unmodified
\texttt{scan\_windows} reported across the five presets, and what fraction of
each window's duration the dense reference showed was actually within OPAS's
own screening radius of a tracked object.

\begin{table}[h]
\centering
\begin{tabular}{lrrr}
\toprule
Preset & Window duration & Reference samples & False-safe rate \\
\midrule
ISS (420~km)        & 24~min & 24 & 100\% \\
Starlink (550~km)   & 46~min & 46 & 100\% \\
Starlink (550~km)   & 16~min & 16 & 100\% \\
SSO (705~km)        & 20~min & 20 & 100\% \\
GEO transfer        & 360~min (clipped) & 360 & 0\% \\
\bottomrule
\end{tabular}
\caption{Baseline false-safe rate for every window OPAS reported as safe,
across a 6-hour scan horizon. Polar orbit reported no windows and is
omitted.}
\label{tab:baseline}
\end{table}

Every window reported in the three crowded low-Earth-orbit shells (ISS,
Starlink, SSO) was, per the dense reference, within OPAS's own danger radius
for the entirety of its reported duration. GEO is the control case: with only
840 tracked objects spread across a far larger volume, OPAS's claim of safety
there was correct.

\subsection{Root cause: coarse sampling, not incorrect physics}

Three properties of the unmodified scanner jointly explain
Table~\ref{tab:baseline}. First, \texttt{scan\_windows}' inner check samples
each orbit at only 12 intervals (about 14 points, roughly one every
6~minutes), against a 50~km screening radius; the dense reference showed that
at this radius, 100\% of one-minute launch times across the scan horizon had
at least one object within range for ISS, Starlink, and SSO alike (mean
8.8--12.2 objects per launch), meaning a real encounter has only a small
chance of landing on one of the sampled points. Second, a parameter sweep
across the coarse step, fine step, and inner-orbit sample count is consistent
with this: doubling the inner sample count from 12 to 24 reduced the reported
window count to zero for all three crowded shells at every coarse step tested,
rather than producing narrower but genuinely safe windows. Third, an
independent boundary-handling defect was identified: the refinement loop that
trims a candidate window's start and end does not clamp to the originally
requested scan range, occasionally extending a reported window several
minutes beyond what was actually evaluated.

\subsection{Engineered fix and its residual limitation}

A corrected scanner was implemented with four changes
(Section~\ref{sec:methodology}): 600-point trajectory sampling (versus
roughly a dozen points), a screening radius of 10--15~km scaled by TLE age
(versus a flat 50~km with an uncapped exponential age penalty reaching up to
200~km), 1-minute temporal verification (versus 2-minute), and clamping of
returned windows to the requested horizon. Applied to the ISS and Starlink
presets, this replaced each preset's single, entirely false-safe window with
one new window per preset (Table~\ref{tab:fixed}); SSO returned no windows,
addressed directly below.

\begin{table}[h]
\centering
\begin{tabular}{llrrr}
\toprule
Preset & Window & Reference samples & Marked unsafe & False-safe rate \\
\midrule
ISS       & 10:33--11:25 (52~min) & 52 & 5 & 9.6\% \\
Starlink  & 10:03--10:22 (19~min) & 19 & 4 & 21.1\% \\
\bottomrule
\end{tabular}
\caption{Corrected scanner's reported windows, validated at the scanner's
exact per-object radius formula (10--15~km by TLE age). SSO returned no
window.}
\label{tab:fixed}
\end{table}

The residual false-safe rate in both windows was examined case by case. For
each of the 10 reference-confirmed object--launch-time pairs within these
two windows' combined 9 unsafe launch minutes (5 for ISS, 4 for Starlink),
the corrected scanner's own 600-point sampled positions were recomputed at
the exact launch instant the reference flagged. All ten (100\%) were
confirmed to be genuinely invisible to the scanner's own math: its sampled
minimum distance at that launch time remained well outside its radius even
though the reference's interpolation-refined minimum was inside it ---
consistent with the scanner's $\sim$9.3--9.6-second waypoint spacing (which
varies slightly with each preset's orbital period) versus LEO relative
velocities of up to $\sim$15~km/s, which can move two objects up to
$\sim$135--145~km between consecutive samples.

\subsection{Direct verification of clear-window existence}

The corrected scanner's coarse-to-refined algorithm examines at most five
candidate coarse-safe runs per scan (Section~\ref{sec:methodology}), and, as
shown above, a window it does report is not guaranteed to be genuinely
clear. To test both possibilities independently of the scanner's own
algorithm, every contiguous stretch of reference-confirmed clear one-minute
launch times was directly enumerated across the full 6-hour horizon at the
corrected radius, for all three crowded shells.

No stretch of at least 15 minutes exists for Starlink (longest 14~min) or
SSO (longest 6~min): no interval of the reported Starlink window's own
length (19~min) is genuinely clear anywhere in the horizon, consistent with
its 21.1\% false-safe measurement in Table~\ref{tab:fixed}, and SSO's
absence of any reported window is correct. ISS, by contrast, contains five
genuinely clear stretches of 15--19~minutes across the horizon; only part of
one overlaps the corrected scanner's single reported window, which is
itself only partially clear (Table~\ref{tab:fixed}). The corrected
scanner's five-run cap is therefore not guaranteed to locate every genuine
clear interval that exists, independent of its residual false-safe rate.

\section{Discussion}
\label{sec:discussion}

\subsection{The original 50~km radius was never meant to be a final threshold}

OPAS's original screening radius (50~km, before TLE-age inflation) is
coincidentally close in magnitude to the 50~km radial extent of one
documented NASA coarse-screening volume~\cite{nasa_screening}. The critical
difference is that NASA's screening volumes are explicitly assigned per
protected asset and orbit regime rather than fixed, and even this large
example was used only as an initial filter within a six-month study, not
applied universally as a final safety threshold; a separate documented
operational example for a near-Earth-orbit asset uses a radius of just
0.5~km~\cite{nasa_screening,nasa_handbook}. OPAS instead applies a single
50~km radius universally, as the sole and final safety criterion,
regardless of orbit regime or mission. This reframes
Section~\ref{sec:results}'s finding: the radius itself was not
unreasonable as one possible coarse filter, but using it as a final,
one-size-fits-all decision threshold --- combined with under-sampling that
never properly evaluated even that generous radius --- produced a tool
that reported safety where none existed.

\subsection{The corrected radius is an assumption, not literature-derived}

The corrected scanner's screening radius (a 10~km base plus 0.5~km per day of
TLE age, capped at 5~km) was chosen for tractability, not derived from
measured SGP4 error rates or from the screening volumes in
Section~\ref{sec:background}. Because the median TLE age in the snapshot is
under one day, the operative radius for most objects is close to the 10~km
base. Its relationship to the literature is not one-directional: the
1--3~km/day heuristic~\cite{jaxsgp4} implies a steeper age penalty than the
0.5~km/day used here, whereas the single-object measurement
in~\cite{afit_sgp4} and the 0.5~km radial extent of a documented NASA
operational volume~\cite{nasa_screening,nasa_handbook} suggest a 10~km sphere is generous in the radial
direction. The Starlink and SSO result is correspondingly radius-sensitive: in
this study's reference data, the fraction of launch times with at least one
object inside a flat radius rises from about 10--11\% at 5~km to 29--45\% at
10~km for those two shells. We therefore present the results in
Section~\ref{sec:results} as a property of the corrected radius, not of any
documented operational threshold, whereas the baseline finding, measured at
OPAS's own radius, does not depend on it.

\subsection{An isotropic radius is itself a simplification}

Both OPAS's original scanner and the corrected version in this study use a
single spherical radius around each tracked object. Documented practice uses an anisotropic ellipsoid instead --- far tighter in the radial direction than along-track or cross-track~\cite{nasa_screening,nasa_handbook} --- because,
as shown in Section~\ref{sec:background}, SGP4/TLE position error is
dominated by in-track drift rather than radial displacement. A sphere large
enough to cover realistic in-track error will therefore be much larger than
necessary in the radial direction, over-restricting some genuinely safe
launch times, while a sphere sized to radial error would under-cover the
in-track direction. This is a structural limitation of OPAS's design, not
something the parameter changes made in this study could resolve, and is
proposed as future work in Section~\ref{sec:conclusion}.

\subsection{The residual false-safe rate has one confirmed cause}

Section~\ref{sec:results} traced the corrected scanner's remaining
9.6--21.1\% false-safe rate, in both windows it reported, entirely to finite
waypoint spacing: every one of the ten reference-confirmed close approaches
examined was invisible to the scanner's own sampled positions at the exact
launch instant the reference flagged, consistent with $\sim$9.3--9.6-second
sample spacing against LEO relative velocities of up to $\sim$15~km/s. No
fixed sampling interval can guarantee catching every possible close approach
at these velocities without approaching the computational cost of
continuous collision detection.

\section{Limitations}
\label{sec:limitations}

\textbf{Scope.} This study evaluates only OPAS's \texttt{/safe-windows}
screening-radius logic (\texttt{scan\_windows} and its inner threat check).
OPAS's separate probability-of-collision scoring, used by its
\texttt{/alert} endpoint, was not evaluated and is outside this study's
scope.

\textbf{Trajectory model.} Both OPAS and this study's reference
implementation model the launch vehicle's path as a simplified circular
orbit derived from launch site, target altitude, and inclination. Real
launch trajectories --- particularly the ascent phase, deliberately
excluded from both OPAS's inner check and this study's reference to keep
the comparison like-for-like --- follow a more complex path shaped by
vehicle performance and guidance constraints. Findings in this report
characterize OPAS's own model, not real-world launch trajectories.

\textbf{Single data snapshot.} All measurements ran against one frozen
catalog snapshot (Section~\ref{sec:methodology}). Debris populations and TLE
freshness vary over time; the specific false-safe rates reported here are a
measurement at one point in time, not a permanent property of OPAS
independent of catalog state. The qualitative finding --- that coarse
sampling can hide close approaches regardless of catalog density --- is
expected to generalize, but the specific percentages in
Table~\ref{tab:baseline} and Table~\ref{tab:fixed} should not be read as
universal constants.

\textbf{Screening radius.} As discussed in Section~\ref{sec:discussion}, the
corrected scanner's 10~km base radius and TLE-age scaling are engineering
choices, not values derived from measured propagation error, and the
Starlink and SSO results are sensitive to them.

\textbf{Candidate-run cap.} \texttt{scan\_windows}, and the corrected
scanner which inherits its structure, examines at most five raw
coarse-safe candidate runs per scan. Section~\ref{sec:results} shows this
can cause a genuinely clear window to go unreported even where one exists
elsewhere in the scanned horizon, independent of the false-safe rate
discussed above.

\textbf{Reference implementation precision.} The dense reference
implementation's launch-time grid is built from an actual floating-point
sample spacing rather than exact clock time, and could in principle drift
from true clock time over a multi-hour scan; no such drift was found to
affect any of the ten residual cases examined in Section~\ref{sec:results},
though this was not independently checked across the reference's full
coverage. Confidence in the reference as ground truth is otherwise
supported by its independent validation (Section~\ref{sec:methodology}).

\section{Conclusion and Future Work}
\label{sec:conclusion}

This study audited OPAS's \texttt{/safe-windows} scanner against an
independently-implemented, validated dense reference, and found that every
window reported as safe in three of five tested orbital regimes was, in
fact, unsafe for its entire reported duration. This was traced to
insufficient trajectory sampling relative to the object density in crowded
low-Earth-orbit shells, not to an error in the underlying orbital
mechanics. A corrected scanner --- denser sampling, a rescaled screening
radius, and finer temporal verification --- eliminated the false-safe
problem's dominant cause. Applied to the three crowded shells, it returned
one window each for ISS and Starlink, with a residual false-safe rate (9.6\%
and 21.1\% respectively) traced entirely to finite waypoint spacing, and no
window for SSO. Direct enumeration of clear intervals in the dense
reference, independent of the scanner's own algorithm, confirmed that
Starlink and SSO genuinely have no clear stretch of at least 15 minutes at
the corrected radius --- so the reported Starlink window is itself
false-safe --- while also showing that the scanner's five-candidate-run cap
can cause it to miss genuinely clear ISS windows that do exist.

Three directions follow directly from the limitations discussed in
Section~\ref{sec:limitations}. First, replacing the isotropic screening
radius with an anisotropic ellipsoid --- tighter radially than along-track,
consistent with documented SGP4 error behavior and real conjunction-screening
practice --- would likely change which launch times are flagged unsafe.
Second, tying the TLE-age radius penalty to measured propagation-error rates,
rather than the tractability-driven formula used here, would ground the
corrected scanner's radius in evidence; its effect on the number of safe
windows cannot be predicted from this study's data alone. Third, the
9-second waypoint gap identified as the corrected scanner's sole confirmed
residual failure mode is a fundamental limit of discrete point-sampling;
closing it fully would require either continuous swept-volume collision
detection or a sampling density approaching the reference implementation's
own computational cost --- a genuine engineering trade-off between accuracy
and the response-time constraints of a live web API, rather than a simple
parameter change.

\begin{thebibliography}{9}

\bibitem{jaxsgp4}
C. Priestley and W. Handley,
``jaxsgp4: GPU-accelerated mega-constellation propagation with batch parallelism,''
\emph{arXiv preprint arXiv:2603.27830}, 2026.

\bibitem{afit_sgp4}
A. T. Rich,
``Investigating Analytical and Numerical Methods to Predict Satellite Orbits Using Two-Line Element Sets,''
thesis, Air Force Institute of Technology, 2017.

\bibitem{nasa_screening}
M. D. Hejduk and D. A. Pachura,
``Conjunction assessment screening volume sizing and event filtering in light of natural conjunction event development behaviors,''
\emph{AAS/AIAA Astrodynamics Specialist Conference}, Stevenson, WA, Aug. 2017.

\bibitem{nasa_handbook}
NASA Office of the Chief Engineer,
``NASA Spacecraft Conjunction Assessment and Collision Avoidance Best Practices Handbook,''
NASA/SP-20230002470, Rev.~1, Feb. 2023.

\bibitem{nasa_cospar}
A. Mashiku,
``NASA spacecraft conjunction assessment and collision avoidance best practices: The history, evolution and opportunities,''
presentation, \emph{44th COSPAR Scientific Assembly}, Athens, Greece, July 2022.

\end{thebibliography}

\end{document}
